\documentclass[14pt,a4paper]{extarticle}
\usepackage[T2A]{fontenc}
\usepackage[utf8]{inputenc}
\usepackage[russian,english]{babel}
\usepackage{lmodern}
\usepackage{mathtext}
\usepackage{geometry}
\sloppy
\geometry{left=30mm,right=30mm,top=20mm,bottom=20mm}
\usepackage{setspace}
\onehalfspacing
\usepackage{indentfirst}
\setlength{\parindent}{1.25cm}
\setlength{\parskip}{0pt}
\usepackage{amsmath,amssymb}
\usepackage{graphicx}
\usepackage{array}
\usepackage{longtable}
\usepackage{booktabs}
\usepackage{caption}
\captionsetup{font=small,labelfont=bf,labelsep=period,justification=centering}
\usepackage{float}
\usepackage{enumitem}
\usepackage{hyperref}
\hypersetup{colorlinks=true,linkcolor=black,urlcolor=blue,citecolor=black,pdfborder={0 0 0}}
\usepackage{titlesec}
\titleformat{\section}{\centering\bfseries\large}{\thesection.}{0.5em}{}
\titleformat{\subsection}{\bfseries\normalsize}{\thesubsection.}{0.5em}{}
\titleformat{\subsubsection}{\bfseries\normalsize}{\thesubsubsection.}{0.5em}{}
\titlespacing*{\section}{0pt}{18pt}{12pt}
\titlespacing*{\subsection}{0pt}{12pt}{6pt}
\titlespacing*{\subsubsection}{0pt}{10pt}{5pt}
\setcounter{secnumdepth}{3}
\setcounter{tocdepth}{2}

% Универсальный плейсхолдер для рисунков.
% При добавлении изображения замените содержимое блока на \includegraphics[width=...]{...}.
\newcommand{\placeholderfigure}[3]{%
  \begin{minipage}[c][#2][c]{#1}
    \centering
    \fbox{\begin{minipage}[c][#2][c]{0.96\linewidth}
      \centering\itshape #3
    \end{minipage}}
  \end{minipage}%
}

\begin{document}

\begin{titlepage}
\thispagestyle{empty}
\begin{center}
Министерство науки и высшего образования Российской Федерации\\
Федеральное государственное автономное образовательное учреждение\\
высшего образования\\
«Национальный исследовательский университет «Московский институт электронной техники»\\[1em]
Институт биомедицинских систем\\[3em]

Пустовит Георгий Максимович\\[3em]

{\Large\bfseries МАГИСТЕРСКАЯ ДИССЕРТАЦИЯ}\\[1em]
по направлению 12.04.04 «Биотехнические системы и технологии»\\[2em]

{\Large\bfseries Восстановление поверхностной электромиограммы высокой плотности из сигналов малоканальных гибких электродов с помощью порождающих моделей}
\end{center}

\vfill

\begin{flushright}
Студент Пустовит Г.М.\\[1em]
Руководитель,\\
к.ф.-м.н., доцент Рябкин Д.И.\\
\end{flushright}

\vfill

\begin{center}Москва 2026\end{center}
\end{titlepage}

\tableofcontents
\clearpage

\section*{Реферат}\addcontentsline{toc}{section}{Реферат}

\section*{Abstract}\addcontentsline{toc}{section}{Abstract}

\section*{Введение}

Электромиография (ЭМГ) -- метод исследования биоэлектрических
потенциалов, возникающих в скелетных мышцах животных и человека при
возбуждении мышечных волокон, электромиограмма -- сигнал, полученный с
помощью метода электромиографии. ЭМГ может быть как инвазивная --
использование игл и микроигл, так и поверхностная -- накожные электроды,
гибкие электроды и высокоплотные электродные матрицы.

В данной работе рассматривается поверхностная ЭМГ, как наиболее
применяемый метод в медицинской практике и устройствах мониторинга
мышечной активности.

Поверхностная ЭМГ применяется для диагностики мышечных заболеваний,
управления бионическими протезами, системах человек-компьютер. Системы с
многоканальной регистрации ограничены усилительной установкой, которая
обладает повышенной стоимостью и габаритами, не позволяющим их
использовать в портативных устройствах.

Для уменьшения числа регистрируемых каналов с сохранением разрешения
сигнала предлагается использовать порождающие модели. Их применение
позволит восстанавливать с контролируемым уровнем ошибки высокоплотный
ЭМГ сигнал, по нескольким каналам, что позволит снизить
энергопотребление датчиков и число каналов усиления, что позволит
сделать шаг на пути к созданию портативных усилителей.

\section{Восстановление высокоплотной электромиограммы по малоканальным сигналам}

Цель работы -- разработать метод и предложить алгоритм восстановления
высокоплотной~поверхностной ЭМГ по малоканальным данным c помощью
порождающих моделей, являющимся вариантом решения компромисса между
точностью регистрации ЭМГ и компактностью регистрирующего устройства.

\subsection{Физиологические основы регистрации и генерации поверхностной электромиографии}

В покое проницаемость мембран клеток намного выше для калия+ чем
натрия+. Поэтому часть ионов калия может выходить из клетки, создавая
снаружи избыток положительных ионов, соответственно внутри избыток
отрицательных ионов. Таким образом создается потенциал мембраны
относительно внутренней части клетки, называемый потенциалом покоя.

У клетки существует два типа каналов: калиевый и натрийевый и по два
типа ворот: активационные и инактивационные. В покое инактивационные
ворота закрыты (рис.1а). Под влиянием раздражителя активационные ворота
натриевого открываются и изменения заряда мембраны что приводит к
изменению потенциала клетки -- потенциалом действия (рис.1б), суммарный
потенциал множества двигательных единиц регистрируется поверхностной
ЭМГ.

Рисунок 1. Изменение потенциала мембраны клетки, а -- иллюстрация
механизма изменения потенциала мембраны клетки, б -- график изменения
потенциала

\begin{figure}[htbp]
\centering
\placeholderfigure{0.42\textwidth}{0.25\textheight}{Изображение 1а: механизм изменения потенциала мембраны}
\hfill
\placeholderfigure{0.42\textwidth}{0.25\textheight}{Изображение 1б: график изменения потенциала}
\caption{Изменение потенциала мембраны клетки.}
\end{figure}

Частотно-амплитудные характеристики сигнала представлены в таблице ниже:

\begin{table}[htbp]
\centering
\caption{Амплитудно-частотные характеристики ЭМГ сигнала.}
\label{tab:emg-characteristics}
\begin{tabular}{p{0.48\textwidth}p{0.32\textwidth}}
\toprule
\textbf{Характеристика} & \textbf{Значение} \\
\midrule
Амплитуда & 0--10~мВ \\
Частота & 0--500~Гц \\
Основной спектр & 50--150~Гц \\
\bottomrule
\end{tabular}
\end{table}

\subsection{Аппаратные средства поверхностной электромиографии}

Датчик поверхностей электромиографии состоит из следующих блоков:
контактной площадки, включающие в себя регистрирующие и опорные
электроды, дифференциальный усилитель, операционный усилитель, фильтр
низких частот, и аналого-цифровой преобразователь, блок схема
представлена на рисунке 2.

\begin{figure}[htbp]
\centering
\placeholderfigure{0.85\textwidth}{0.22\textheight}{Плейсхолдер: схема устройства датчика поверхностной электромиографии с системой предобработки}
\caption{Схема устройства датчика поверхностной электромиографии с системой предобработки.}
\end{figure}

Для ЭМГ датчиков коэффициент ослабления синфазного сигнала должен
составлять порядка 60-100 дБ, данной характеристикой определяется выбор
дифференциального усилителя:

\begin{equation}
\text{КОСС}\,(\text{дБ}) = 20\log_{10}\left(\frac{V_{\text{вых}}}{V_{\text{сф}}}\right).
\end{equation}

где V\textsubscript{сф}, -- напряжение синфазного сигнала,
V\textsubscript{вых} -- выходное напряжение

\subsubsection{Биполярные датчики: типы электродов, схемы включения, требования к коэффициенту ослабления синфазного сигнала}

\subsubsection{Гибкие электродные системы}

\subsubsection{Электродные матрицы высокой плотности и ограничения масштабирования числа каналов}

\subsection{Методы восстановления высокоплотной электромиограммы по неполным данным}

Описание принципа работы и обзор применения для каждого из методоа для
восстановления ЭМГ сигналов. Ограничения каждого метода.

\subsubsection{Классические методы: пространственная интерполяция, методы алгоритмы машинного обучения}

\subsubsection{Свёрточные нейронные сети}

\subsubsection{Порождающие модели}

\subsection{Выводы}

\section{Метод восстановления электромиограммы на основое маскированного автокодировщика}

\subsection{Формальная постановка задачи}

\subsection{Предобработка сигналов: фильтрация, разбиения на окна, экстраполяция}

\subsection{Архитектура маскированного автокодировщика и обоснование ее выбора}

\placeholderfigure{0.82\textwidth}{0.25\textheight}{Плейсхолдер: архитектура маскированного автокодировщика}
\begin{center}\textit{Замените этот блок схемой архитектуры модели.}\end{center}

\subsection{Стратегии маскирования и разрыв между обучающими и рабочими масками}

\subsection{Метрики качества предсказания модели и вычислительной сложности моделей}

\subsection{Выводы}

\section{Реализация и обучение порождающей модели}

\subsection{Описание набора данных и дизайна эксперимента}

\placeholderfigure{0.82\textwidth}{0.24\textheight}{Плейсхолдер: схема экспериментальной установки и расположения электродов}

Таблица с описанием наборов + описание эксперимента и оборудования

\subsection{Программная реализация и настройка обучения}

\subsection{Сравнение стратегий маскирования}

\placeholderfigure{0.82\textwidth}{0.24\textheight}{Плейсхолдер: график сравнения стратегий маскирования}

\subsection{Влияние размера фрагмента на качество восстановления}

\subsection{Сравнение с базовыми методами с GenENet}

\placeholderfigure{0.82\textwidth}{0.24\textheight}{Плейсхолдер: сравнение с базовыми методами и GenENet}

\subsection{Ограничения}

Выводы

\end{document}
